% Insert after the mixed-jet equation in the node calculation.
The displayed condition is also sufficient in the analytic local ring.
Write $R=\mathbb C\{x,y\}$ and let $\varphi:R^4\to R$ be the
actual cup row, whose image is $(x,y)$. A pair of branch kernel sections
with a common value has a lift $\widetilde\eta\in R^4$ satisfying
$\varphi(\widetilde\eta)=xyh$. The constant term of $h$ is exactly
$\langle H,\eta_t'(0)\rangle-\langle B,\eta_h'(0)\rangle$.
Adding $xy\zeta$ to the lift changes $h$ by $\varphi(\zeta)$, and
these changes range over the whole ideal $(x,y)$. Once the displayed
constant vanishes, choose $\zeta$ with $\varphi(\zeta)=-h$.
This gives a section of the total-space kernel with precisely the
prescribed branch restrictions. Thus the one mixed-jet equation is the
complete remaining node condition, including all higher analytic terms.

% Insert immediately before the calculation using H=I and diagonal B.
Here is the actual change of frames used in that calculation.
For new frames $F'=FU$ and $G'=GV$, the derivative and tensor matrices
transform by
\[
 T'=V^{-1}TU,\qquad \eta'=U^{-1}\eta V^{-t},\qquad
 \langle T',\eta'\rangle=(\det V)^{-1}\langle T,\eta\rangle.
\]
The last equality follows from
$V^{-t}JV^{-1}=(\det V)^{-1}J$, where
$J=\left(\begin{smallmatrix}0&1\\-1&0\end{smallmatrix}\right)$.
The conic matrices become $A'=U^{-1}AU$ and $B'=V^{-1}BU$.
Choose an eigenvector matrix $R$ for the actual endomorphism $H^{-1}B$,
and set $U=R$ and $V=HR$. These maps send $H$ to $I$, send $B$ to
$\operatorname{diag}(\beta_1,\beta_2)$, and send the tensor to
$R^{-1}\eta H^{-t}R^{-t}$. They preserve the kernel equations through
the displayed determinant-line isomorphism and give the diagonal
calculation in the declared frames.
